\chapter{Questão de Pesquisa}

Como implementar e otimizar algoritmos de visão computacional e modelos de Deep Learning para identificação e localização de objetos em tempo real, operando estritamente dentro das restrições de processamento de sistemas embarcados (SBCs como Raspberry Pi e Microcontroladores ESP32), garantindo autonomia ao braço robótico Atom sem depender de computadores externos?

\subsection{Justificativa}
A pesquisa justifica-se pela necessidade de democratizar o acesso à robótica inteligente. Validar uma arquitetura que permita rodar IA avançada em hardware acessível (como o ESP32-CAM e Raspberry Pi) não apenas viabiliza o projeto Atom, mas contribui para o campo da Robótica de Serviço e IoT (Internet of Things). Além disso, o estudo aprofunda a análise técnica iniciada anteriormente sobre a eficiência de software, movendo-se da teoria de comparação de linguagens (Python vs. C++)  para a aplicação prática em engenharia de sistemas embarcados. A autonomia do sistema elimina a necessidade de tethering (conexão cabeada a um PC), aumentando a versatilidade do robô.